\section{The finite Paley--Wiener trace and the order of its limits}
\label{sec:global-pw-trace}

This section uses the actual range of MA-GP-003 and the specified
symmetric weight.  In particular it retains the bounded Gram change
of MA-GP-004.  It proves trace class of the actual convolution after
the orthogonal projection, before computing its trace.
All test functions in this section are additive tests
$g\in C_c^\infty((0,\infty))$, extended by zero to $\mathbb R$.
This support condition makes the half-line subspace and its
finite-level ideal invariant under convolution.  In multiplicative
coordinates the corresponding test is supported in $(1,\infty)$.

\paragraph{MA-GP-009 (the actual finite-level trace and two iterated limits).}
Let $0<\alpha<1/2$, $s=0$, and
$H=L^2(\mathbb R,e^{2\alpha|x|}\,dx)$.
For $Y>0$ set $c=-\log Y$ and
\[
 H_c=\{f\in H:\operatorname{supp}f\subseteq[c,\infty)\},
 \qquad \mathcal Q_c=H_c/J_Y,\qquad Z_c=H_c\ominus J_Y .
 \tag{MA-GP.80}
\]
Here $J_Y$ and $Q_Y^{\rm glob}=P_{J_Y}$ are the actual
Paley--Wiener closure and orthogonal projection of MA-GP-001--003.
Let $P_c^Z$ denote the orthogonal projection onto $Z_c$,
extended by zero on $H_c^\perp$.
Then the actual operators satisfy
\[
 \begin{gathered}
 C_gP_c^Z\in\mathcal S_1(H),\qquad
 C_g^{\mathcal Q_c}\in\mathcal S_1(\mathcal Q_c),\\
 \operatorname{Tr}_H(C_gP_c^Z)
 =\operatorname{Tr}_{\mathcal Q_c}(C_g^{\mathcal Q_c})
 =\sum_{\operatorname{Re}\rho<1/2+\alpha}
           m_\rho\,\mathcal Mg(\rho-1/2).
 \end{gathered}
 \tag{MA-GP.81}
\]
The sum is over nontrivial zeta zeros, with multiplicity, and
converges absolutely.  It is independent of $Y$.

Choose the uniformly bounded smooth compact exhaustion
$\chi_R(x)=\chi_0(x/R)$ with $\chi_0=1$ near zero.
Both iterated limits of the actual localized operators exist:
\[
 \begin{aligned}
 \lim_{Y\to\infty}\lim_{R\to\infty}
   \operatorname{Tr}(M_{\chi_R}C_gQ_Y^{\rm glob})
   &=-\sum_{\operatorname{Re}\rho<1/2+\alpha}
                  m_\rho\,\mathcal Mg(\rho-1/2),\\
 \lim_{R\to\infty}\lim_{Y\to\infty}
   \operatorname{Tr}(M_{\chi_R}C_gQ_Y^{\rm glob})
   &=-\sum_{1/2-\alpha<\operatorname{Re}\rho<1/2+\alpha}
                  m_\rho\,\mathcal Mg(\rho-1/2).
 \end{aligned}
 \tag{MA-GP.82}
\]
Their difference, in the displayed order, is
\[
          -\sum_{\operatorname{Re}\rho\le1/2-\alpha}
                         m_\rho\,\mathcal Mg(\rho-1/2).
 \tag{MA-GP.83}
\]
The formula includes the lower-edge zeros.  It does not assert
that such exterior or boundary zeros occur, or that the
difference is nonzero for any specified test.

\begin{proof}
\emph{The actual quotient action and its Hardy representative.}
Put
\[
 (U_cf)(y)=e^{\alpha(c+y)}f(c+y),\qquad
 R_c(y)=e^{\,2\alpha|c+y|-2\alpha(c+y)},\qquad y\ge0.
 \tag{MA-GP.84}
\]
The map $U_c$ is unitary from the actual $H_c$ onto
$H_{R_c}=L^2((0,\infty),R_c(y)\,dy)$.
For this fixed $c$ the norm is equivalent to ordinary $L^2$,
since $1\le R_c\le e^{4\alpha\max(-c,0)}$.
If $c\ge0$ it is identically one; if $c<0$ and $L=-c$,
\[
 R_c(y)=
 \begin{cases}
       e^{4\alpha(L-y)},&0\le y\le L,\\
       1,&y\ge L.
 \end{cases}
 \tag{MA-GP.85}
\]
MA-GP-003 says that, as an actual closed set,
the image of $J_Y$ under the Laplace map $\mathcal L U_c$
is $B_\alpha H^2(\operatorname{Re}z>0)$.
Here $B_\alpha$ is the pure Blaschke product with zeros
\[
                 z_j=\alpha+\tfrac12-\rho_j=a_j+ib_j,
          \qquad a_j>0,
 \tag{MA-GP.86}
\]
repeated according to multiplicity.
The norms induce a bounded invertible map between their
quotient Hilbert spaces.  Thus the actual quotient action is
boundedly similar to the standard model-space action on
$K_{B_\alpha}=H^2\ominus B_\alpha H^2$.
This statement concerns the action on the quotient; it does
not replace the actual weighted orthogonal projection by the
ordinary Hardy projection.

The unilateral shifts $S_t v(y)=\mathbf1_{y\ge t}v(y-t)$,
$t\ge0$, are contractions on $H_{R_c}$ because $R_c$ is
nonincreasing.  They preserve the image of $J_Y$.
For
\[
      k(t)=e^{\alpha t}g(t),\qquad
      G(z)=\int_0^\infty k(t)e^{-zt}\,dt
                           =\mathcal Mg(\alpha-z),
 \tag{MA-GP.87}
\]
the map $U_c C_g U_c^{-1}$ on $H_{R_c}$ is the actual
causal convolution $C_k^+=\int_0^\infty k(t)S_t\,dt$.
The vectorwise integrals define bounded operators.
The corresponding Hardy quotient action is
\[
                 T_G=P_{K_{B_\alpha}}M_G|_{K_{B_\alpha}}.
 \tag{MA-GP.88}
\]
Since $k$ is smooth and compactly supported inside $(0,\infty)$,
integration by parts gives, for every $N$,
$|G(z)|\le C_N(1+|z|)^{-N}$ on $\operatorname{Re}z\ge0$.
In particular all multipliers used below are bounded analytic
multipliers of $H^2$.

\emph{Uniform half-plane estimates at arbitrarily close edge zeros.}
Use the boundary norm $(2\pi)^{-1}\int|F(it)|^2dt$.
The normalized Malmquist basis of $K_{B_\alpha}$ is
\[
 e_j(z)=\frac{\sqrt{2a_j}}{z+\overline{z_j}}
                         \prod_{\ell<j} b_\ell(z),
 \qquad |b_\ell(it)|=1,\quad b_\ell(z_\ell)=0.
 \tag{MA-GP.89}
\]
Unimodular normalization constants in $b_\ell$ do not affect
the argument.  The elementary orthogonal decomposition
$K_{UV}=K_U\oplus U K_V$ for inner $U,V$ proves orthonormality,
including repeated zeros.  It also proves completeness.
Indeed the finite-product subspaces increase, and their
orthogonal complements decrease to $B_\alpha H^2$:
if $F$ is divisible in $H^2$ by every finite product $B_N$,
then $F/B_N$ has bounded $H^2$ norm.  A weakly convergent
subsequence, evaluated at each point in the half-plane,
has limit $F/B_\alpha$ after the standard normalizations
of $B_N$.  Hence $F/B_\alpha\in H^2$.
This proves the claimed intersection and the basis assertion.
The empty or finite product cases have the same interpretation.

For a bounded analytic multiplier $F$ its basis-vector norm is
\[
 \|M_F e_j\|_{H^2}^2
   =\frac1\pi\int_{\mathbb R}
         |F(it)|^2\frac{a_j}{(t-b_j)^2+a_j^2}\,dt .
 \tag{MA-GP.90}
\]
Suppose $|F(it)|\le C/(1+|t|)$.
There is a constant $C_F'$ independent of $j$ such that
\[
                   \|M_F e_j\|_{H^2}^2
                              \le\frac{C_F'}{1+b_j^2}.
 \tag{MA-GP.91}
\]
To prove it, the containing strip for zeta gives
$0<a_j\le\alpha+1/2<1$.
For $|b_j|\ge2$, split the integral into
$|t-b_j|\le|b_j|/2$ and its complement.
On the first part $|t|\ge|b_j|/2$, and the Poisson
kernel has total mass at most one.
On the second part its density is at most
$4(\alpha+1/2)/(\pi b_j^2)$; use
$\int|F(it)|^2dt<\infty$.
These give (MA-GP.91); bounded $|b_j|$ follows directly
from the total Poisson mass and boundedness of $F$.
The unconditional local zero count NT-Z10 yields
\[
                      \sum_j(1+b_j^2)^{-1}<\infty.
 \tag{MA-GP.92}
\]
No positive lower bound on $a_j$ has been used.
It follows from (MA-GP.89)--(MA-GP.92) that
$M_F|_{K_{B_\alpha}}$ and its compression $T_F$ are
Hilbert--Schmidt.

Fix $\kappa>0$.  Apply this result to
\[
                    r(z)=(\kappa+z)^{-1},
               \qquad F_2(z)=(\kappa+z)G(z).
 \tag{MA-GP.93}
\]
Both meet the bounded-analytic and boundary-decay requirements.
The quotient compressions form a homomorphism:
$T_FT_H=T_{FH}$, because $B_\alpha H^2$ is invariant under
every bounded analytic multiplier.
Consequently
\[
                          T_G=T_rT_{F_2}\in\mathcal S_1,
             \qquad T_r^2\in\mathcal S_1.
 \tag{MA-GP.94}
\]
This proves quotient trace class before a trace is assigned.
Trace in the actual orthonormal basis (MA-GP.89) and the
analytic Poisson formula now give
\[
 \begin{split}
 \operatorname{Tr}T_G
   &=\sum_j\frac1\pi\int_{\mathbb R}
                   G(it)\frac{a_j}{(t-b_j)^2+a_j^2}\,dt\\
   &=\sum_jG(z_j)
    =\sum_{\operatorname{Re}\rho<1/2+\alpha}
                      m_\rho\,\mathcal Mg(\rho-1/2).
 \end{split}
 \tag{MA-GP.95}
\]
The Poisson identity follows also by the Cauchy integral
formula on the right half-plane, using the rapid decay of
$G$.  Absolute convergence follows both from trace class
for the diagonal sum and directly from rapid vertical
decay of $\mathcal Mg$ on the fixed containing strip and
NT-Z10.  Bounded similarity proves the quotient trace
assertion of (MA-GP.81).

\emph{A trace-class resolvent square on the actual adjoint subspace.}
The generator $A_c$ of the quotient semigroup induced by $S_t$
has resolvent represented by $T_r$ in the Hardy quotient.
Indeed the vectorwise strong Laplace integral
$\int_0^\infty e^{-\kappa t}S_t\,dt$ is bounded, and
its multiplier is $(\kappa+z)^{-1}$.
The usual resolvent identity follows directly by changing
variables in its shift difference quotient:
$A_cR_\kappa=\kappa R_\kappa-1$; integration by parts
on the generator domain gives the inverse identity there.
Its tail has operator norm at most $e^{-\kappa T}/\kappa$,
although no operator-norm continuity of $t\mapsto S_t$
is required.  By (MA-GP.94) and bounded similarity,
\[
                        (\kappa-A_c)^{-2}\in\mathcal S_1.
 \tag{MA-GP.96}
\]

Write $R=R_c$ in the remainder of this calculation, and
let $Z_R=U_cZ_c$, with its actual $H_R$ orthogonal projection
$P_{Z_R}$.  The adjoint unilateral semigroup is
\[
  (S_t^{*,R}u)(y)=\frac{R(y+t)}{R(y)}u(y+t),
       \qquad \mathcal L_Ru=R^{-1}\partial_y(Ru).
 \tag{MA-GP.97}
\]
It preserves $Z_R$ because the ideal is invariant under $S_t$.
Multiplication by $R$ is unitary from $H_R$ to
$L^2((0,\infty),R^{-1}dy)$ and converts this adjoint
semigroup to left translation $v(y)\mapsto v(y+t)$.
Its generator has no boundary condition at zero.
Precisely, for $j\le2$ the graph domain consists of
$v=Ru$ with weak derivatives through order $j$ in this
weighted $L^2$, and $\mathcal L_R^j u=R^{-1}v^{(j)}$.
For the first derivative the direct assertion follows from
the fundamental theorem of calculus and strong translation
continuity; the converse follows by testing the generator
difference quotient against interior smooth tests.
Iteration proves the second-derivative statement.
The weight is bounded above and below, so $v\in H^2(0,\infty)$
when $u\in\mathcal D(\mathcal L_R^2)$.

The restriction of (MA-GP.97) to $Z_R$ is the Hilbert
adjoint of the actual quotient semigroup.  Equation
(MA-GP.96), adjunction, and the restricted resolvent integral
therefore give
\[
 \mathcal R_2=(\kappa-\mathcal L_R)^{-2}P_{Z_R}
                       \in\mathcal S_1(H_R).
 \tag{MA-GP.98}
\]
It maps boundedly into the second graph domain, since
$\mathcal L_R(\kappa-\mathcal L_R)^{-1}
=\kappa(\kappa-\mathcal L_R)^{-1}-1$.

\emph{The actual half-line lift, including its endpoint and cusp.}
Set $a(y)=R(y)^{-1}$.  For
$u\in\mathcal D(\mathcal L_R^2)$ and $v=Ru$,
\[
 C_k^+(\kappa-\mathcal L_R)^2u(x)
       =\int_0^\infty k(x-y)a(y)
                             (\kappa-\partial_y)^2v(y)\,dy .
 \tag{MA-GP.99}
\]
The test $k$ is extended by zero on the negative half-line.
Integrate twice by parts on $(0,L)$ and $(L,\infty)$
when $c<0$, or on $(0,\infty)$ when $c\ge0$.
All terms at infinity vanish because $k(x-y)$ has compact
support in $y$.  The bulk operator is
\[
 \mathcal B_k u(x)=
   \int_{y>0,\ y\ne L} R(y)
       (\kappa+\partial_y)^2
                 \bigl(k(x-y)a(y)\bigr)u(y)\,dy ,
 \tag{MA-GP.100}
\]
where the exclusion of $L$ is omitted when $c\ge0$.
After expansion its kernel is the sum of
\[
  [\kappa^2k-2\kappa k'+k''](x-y)
 +[2\kappa k-2k'](x-y)R(y)a'(y)
 +k(x-y)R(y)a''(y).
 \tag{MA-GP.101}
\]
The piecewise functions $Ra'$ and $Ra''$ are bounded;
they equal $4\alpha,16\alpha^2$ on $(0,L)$ and zero
outside when $c<0$.  Hence $\mathcal B_k$ is bounded
on $H_R$, as a finite sum of bounded causal convolutions
and bounded multiplications.

The exact endpoint contribution at zero is
\[
 \bigl[a(0)(2\kappa k-k')+a'(0+)k\bigr]v(0)
                                  -a(0)k\,v'(0).
 \tag{MA-GP.102}
\]
When $c<0$ the only interior contribution is
\[
          [a']_{L-}^{L+}\,k(\,\cdot-L)v(L)
                      =-4\alpha\,k(\,\cdot-L)v(L).
 \tag{MA-GP.103}
\]
To verify the signs, the boundary expression before summing
the intervals is
$qv'-(2\kappa q+q')v$, where $q(y)=k(x-y)a(y)$.
Subtracting its value at zero gives (MA-GP.102).
At $L$, the continuity of $a,v,v'$ cancels all terms
except the jump of $a'$ times $k(x-L)v(L)$.
For $R=1$ these formulas reduce to the explicit identity
\[
 C_k^+(\kappa-\partial)^2u
  =C_{\kappa^2k-2\kappa k'+k''}^+u
                 +(2\kappa k-k')u(0)-k\,u'(0).
 \tag{MA-GP.104}
\]
Thus the half-line endpoint has not been dropped.

The functions multiplying $v(0),v'(0),v(L)$ in
(MA-GP.102)--(MA-GP.103) belong to $H_R$.
The corresponding evaluation functionals are bounded
on the graph domain of $\mathcal L_R^2$.
This follows from the one-dimensional Sobolev trace
inequality on a fixed interval containing the relevant
endpoint, and norm equivalence there.  One can prove
the inequality by averaging the fundamental theorem
of calculus and applying Cauchy--Schwarz to $v$ and
$v'$, then to $v'$ and $v''$.
Consequently their compositions with $\mathcal R_2$
are bounded functionals on $H_R$.

Apply (MA-GP.99)--(MA-GP.103) to
$u=\mathcal R_2f$.  Its left side is $C_k^+P_{Z_R}f$.
The right side is the product of the bounded
$\mathcal B_k$ with the trace-class $\mathcal R_2$,
plus the displayed finite-rank boundary terms.
This proves
\[
                         C_k^+P_{Z_R}\in\mathcal S_1(H_R).
 \tag{MA-GP.105}
\]
Conjugating by the actual unitary $U_c$, and extending by
zero on $H_c^\perp$, proves $C_gP_c^Z\in\mathcal S_1(H)$.
The support of $g$ ensures that $C_g$ maps $H_c$ into itself.
Trace cyclicity with the actual orthogonal projection gives
\[
 \operatorname{Tr}(C_gP_c^Z)
       =\operatorname{Tr}(P_c^ZC_gP_c^Z)
       =\operatorname{Tr}_{\mathcal Q_c}(C_g^{\mathcal Q_c}).
 \tag{MA-GP.106}
\]
The middle compression is the orthogonal representative
of quotient convolution.  Together with (MA-GP.95),
this finishes (MA-GP.81) by an actual lift proof.

\emph{Position cutoff first and the comparison of orders.}
Let $P_c$ be multiplication by $\mathbf1_{[c,\infty)}$.
As actual orthogonal projections on $H$,
\[
                        Q_Y^{\rm glob}=P_c-P_c^Z.
 \tag{MA-GP.107}
\]
For compactly supported smooth $\chi$, the operator
$M_\chi C_gP_c$ is trace class and has trace zero.
Indeed $M_\chi C_g$ is trace class by the actual compact
kernel argument of MA-LIFT-003, and right multiplication
by $P_c$ preserves trace class.
Approximate $P_c$ strongly by uniformly bounded smooth
multiplication functions.  The corresponding products
converge in trace norm, by the finite-rank approximation
argument and adjoint strong convergence.
Each approximating trace is the ordinary diagonal integral
$g(0)\int\chi(x)\eta(x)\,dx=0$.
The weight is handled by the same two bounded compact
coordinate multiplication factors as in MA-LIFT-003;
no smoothness of the weight at zero is required.
Thus
\[
 \operatorname{Tr}(M_{\chi_R}C_gQ_Y^{\rm glob})
                   =-\operatorname{Tr}(M_{\chi_R}C_gP_c^Z).
 \tag{MA-GP.108}
\]
For fixed $Y$, the uniformly bounded multipliers
$M_{\chi_R}$ converge strongly to one.
Since $C_gP_c^Z$ is trace class, this proves
\[
          \|(M_{\chi_R}-1)C_gP_c^Z\|_1\longrightarrow0.
 \tag{MA-GP.109}
\]
Equations (MA-GP.108)--(MA-GP.109) and (MA-GP.81)
give the first limit in (MA-GP.82); its inner limit
is already independent of $Y$.

For the opposite order, MA-GP-001 gives the actual
fixed-cutoff trace-norm convergence as $Y\to\infty$,
and then the all-weight ambient theorem MA-WP-003
gives the full-strip quotient trace.
At $s=0$ no boundary jet is continuous, since the
edge condition is $2j<s-1=-1$.
Thus precisely the strict strip
$1/2-\alpha<\operatorname{Re}\rho<1/2+\alpha$
is retained.  This is the second limit in (MA-GP.82).
Subtracting the two absolutely convergent sums proves
(MA-GP.83), including the lower-edge convention.
\end{proof}

\paragraph{Scope, source placement, and the remaining transport problem.}
All estimates above hold for a fixed finite $Y$.
The bound relating the actual symmetric weight to the half-line
exponential weight grows as $Y^{4\alpha}$ for $Y>1$.
No uniform trace-norm bound for the lifted operators as
$Y\to\infty$, no joint-limit theorem, and no prescribed
diagonal rate are inferred.
The defect (MA-GP.83) is an exact statement about two orders
of limits of the actual global Paley--Wiener projections.
It is separate from the local additive Fourier-cutoff
calculation in MA-LA-014.

The proof relies on the actual range theorem MA-GP-003,
with its primary Beurling/Burnol locators, and on NT-Z10's
unconditional zero count.  The half-plane Malmquist calculation
is supplied explicitly; it is the half-plane version of the
compressed-multiplier argument in MA-HY-002.
The actual lift is a half-line version of MA-LIFT-002 with
the endpoint at zero and the weight cusp both retained.
These are proofs for the specified Hilbert model, not an
appeal to the source's uncorrected ambient nuclearity claim.
Neither $\mathsf A$ nor $M$ is used.

In the inherited Grothendieck trace framework, the finite
zero jets give the finite-length spectral blocks and the
bounded-to-trace-class factorization makes the ambient
trace legitimate before it is evaluated.  The applicable
cote~1 locators are the already recorded scans~079--080
and~111--112, ALPOGE\_BRIDGES BR-TR-01/02.
No additional archival reading or new arithmetic motive
realization is asserted by this operator calculation.
